% Prose rendering of PrimeTensor/Prime.lean.
% Fragment: no \documentclass, no preamble, no document environment.

\section{Prime atoms, prime bags, and prime multisets}
\label{Prime:sec:prime}

\leanfile{PrimeTensor/Prime.lean}

\subsection*{What this file establishes}

The file builds the multiplicative carrier the project starts from, in two
stages.  First an \emph{ordered} factor chain, $\name{PrimeBag}$, whose
terminator is the multiplicative pivot $\ctor{one}$ and whose multiplication
is concatenation.  Concatenation is associative and unital but visibly not
commutative: the chain $2 \cdot 3$ and the chain $3 \cdot 2$ are different
terms.  Second, the quotient $\name{PrimeMultiset}$ of chains by equality of
represented magnitude, on which multiplication \emph{is} commutative --- not
up to a relation, but as an equality of elements.

So the file's arc is: pick a representation in which the operation is
computable and order-sensitive, then quotient exactly enough that order stops
being observable at the public boundary.

\subsection{Prime atoms}

\begin{definition}[Prime, \lean{PrimeTensor.Prime}]\label{Prime:def:prime}
$\name{Prime}$ is the structure with fields
\[
  \name{value} : \mathbb{N},
  \qquad
  \name{prime} : \name{Nat.Prime}(\name{value}).
\]
\end{definition}

\begin{lstlisting}
structure Prime where
  value : ℕ
  prime : Nat.Prime value
\end{lstlisting}

An atom is a natural number bundled with a proof that it is prime, so
primality travels with the number and never has to be re-established.  Two
atoms are equal exactly when their $\name{value}$ fields are --- the proof
field is a proposition, hence proof-irrelevant.

\begin{lemma}[\lean{PrimeTensor.Prime.one\_lt}]\label{Prime:lem:one-lt}
For every $p : \name{Prime}$, $\; 1 < p.\name{value}$.
\end{lemma}

\begin{proof}
This is \lean{Nat.Prime.one\_lt} applied to the stored field
$p.\name{prime}$.  It is tagged \lean{@[simp]} so that the strict lower bound
on an atom is available to the simplifier without being re-derived.
\end{proof}

\subsection{Prime bags}

\begin{definition}[Prime bag, \lean{PrimeTensor.PrimeBag}]\label{Prime:def:primebag}
$\name{PrimeBag}$ is the inductive type generated by
\[
  \ctor{one} : \name{PrimeBag},
  \qquad
  \ctor{factor} : \name{Prime} \to \name{PrimeBag} \to \name{PrimeBag}.
\]
\end{definition}

\begin{lstlisting}
inductive PrimeBag where
  | one : PrimeBag
  | factor : Prime → PrimeBag → PrimeBag
\end{lstlisting}

A bag is a finite list of atoms terminated by the pivot; the name records the
intent, but at this stage the structure is a list and the order of factors is
part of the term.  As with $\name{Depth}$ in
\texttt{PrimeTensor/Depth.lean}, the terminator is $\ctor{one}$, the multiplicative
unit, and not a zero.

\begin{definition}[Evaluation, \lean{PrimeTensor.PrimeBag.eval}]
\label{Prime:def:bag-eval}
$\name{eval} : \name{PrimeBag} \to \mathbb{N}$ is defined by structural
recursion:
\begin{align}
  \name{eval}(\ctor{one}) &= 1, \label{Prime:eq:eval-one}\\
  \name{eval}(\ctor{factor}\;p\;r) &= p.\name{value} \cdot \name{eval}(r).
    \label{Prime:eq:eval-factor}
\end{align}
\end{definition}

\begin{lstlisting}
def eval : PrimeBag → ℕ
  | .one => 1
  | .factor p rest => p.value * rest.eval
\end{lstlisting}

\begin{lemma}[\lean{PrimeBag.eval\_one}, \lean{PrimeBag.eval\_factor}]
\label{Prime:lem:eval-eqs}
Equations \eqref{Prime:eq:eval-one} and \eqref{Prime:eq:eval-factor} hold.
\end{lemma}

\begin{proof}
Each is the defining equation of $\name{eval}$ at a constructor, so both sides
are definitionally equal and the proof is $\ctor{rfl}$.  Both are tagged
\lean{@[simp]}: the statements are content-free, and the tag is what makes
them usable by the simplifier, which does not unfold recursive definitions on
its own.
\end{proof}

\begin{lemma}[\lean{PrimeTensor.PrimeBag.one\_le\_eval}]\label{Prime:lem:one-le-eval}
For every $b : \name{PrimeBag}$, $\; 1 \leq \name{eval}(b)$.
\end{lemma}

\begin{proof}
By structural induction on $b$.

For $b = \ctor{one}$ the goal is $1 \leq 1$, closed by reflexivity of $\leq$.

For $b = \ctor{factor}\;p\;r$, Lemma~\ref{Prime:lem:one-lt} gives
$1 < p.\name{value}$ and hence $1 \leq p.\name{value}$, while the induction
hypothesis gives $1 \leq \name{eval}(r)$.  Monotonicity of multiplication on
$\mathbb{N}$ in both arguments then yields
\[
  1 = 1 \cdot 1 \;\leq\; p.\name{value} \cdot \name{eval}(r)
    = \name{eval}(\ctor{factor}\;p\;r). \qedhere
\]
\end{proof}

So evaluation lands in $\mathbb{Z}_{\geq 1}$: no bag represents $0$.  This is
the same decision one layer up from the absence of a zero depth in
\texttt{PrimeTensor/Depth.lean} --- the carrier is positive by construction,
so the multiplicative structure below never meets an absorbing element.

\subsection{Multiplication is concatenation}

\begin{definition}[Product of bags, \lean{PrimeTensor.PrimeBag.mul}]
\label{Prime:def:bag-mul}
$\name{mul} : \name{PrimeBag} \to \name{PrimeBag} \to \name{PrimeBag}$ is
defined by structural recursion on its first argument:
\begin{align}
  \ctor{one} \cdot b &= b, \label{Prime:eq:mul-one-left}\\
  (\ctor{factor}\;p\;r) \cdot b &= \ctor{factor}\;p\;(r \cdot b).
    \label{Prime:eq:mul-factor}
\end{align}
The instances \lean{One PrimeBag} and \lean{Mul PrimeBag} make $1$ denote
$\ctor{one}$ and $\cdot$ denote $\name{mul}$.
\end{definition}

\begin{lstlisting}
def mul : PrimeBag → PrimeBag → PrimeBag
  | .one, b => b
  | .factor p rest, b => .factor p (mul rest b)

instance : One PrimeBag := ⟨.one⟩
instance : Mul PrimeBag := ⟨mul⟩
\end{lstlisting}

This is list concatenation: the factors of $a$ are re-emitted in order and the
chain is then continued by $b$.

\begin{lemma}[Left unit, \lean{PrimeTensor.PrimeBag.one\_mul}]
\label{Prime:lem:bag-one-mul}
$1 \cdot b = b$ for every $b : \name{PrimeBag}$.
\end{lemma}

\begin{proof}
This is \eqref{Prime:eq:mul-one-left}, the defining equation at $\ctor{one}$.  The
recursion matches on the first argument, which is already a constructor here,
so the two sides are definitionally equal and the proof is $\ctor{rfl}$.
\end{proof}

\begin{lemma}[Right unit, \lean{PrimeTensor.PrimeBag.mul\_one}]
\label{Prime:lem:bag-mul-one}
$b \cdot 1 = b$ for every $b : \name{PrimeBag}$.
\end{lemma}

\begin{proof}
By structural induction on $b$; unlike Lemma~\ref{Prime:lem:bag-one-mul} this is
\emph{not} definitional, because the recursion inspects the first argument
and $b$ is a variable.

For $b = \ctor{one}$, \eqref{Prime:eq:mul-one-left} gives
$\ctor{one} \cdot 1 = 1 = \ctor{one}$, again by reflexivity.

For $b = \ctor{factor}\;p\;r$, \eqref{Prime:eq:mul-factor} rewrites the goal to
$\ctor{factor}\;p\;(r \cdot 1) = \ctor{factor}\;p\;r$, and the induction
hypothesis $r \cdot 1 = r$ gives it after congruence under
$\ctor{factor}\;p\;(-)$.
\end{proof}

\begin{lemma}[Associativity, \lean{PrimeTensor.PrimeBag.mul\_assoc}]
\label{Prime:lem:bag-assoc}
$(a \cdot b) \cdot c = a \cdot (b \cdot c)$ for all
$a, b, c : \name{PrimeBag}$.
\end{lemma}

\begin{proof}
By structural induction on $a$, with $b$ and $c$ fixed.

For $a = \ctor{one}$ both sides reduce to $b \cdot c$ by
\eqref{Prime:eq:mul-one-left}, so the case is definitional.

For $a = \ctor{factor}\;p\;r$, two applications of \eqref{Prime:eq:mul-factor} on
each side turn the goal into
$\ctor{factor}\;p\;((r \cdot b) \cdot c)
  = \ctor{factor}\;p\;(r \cdot (b \cdot c))$,
which follows from the induction hypothesis by congruence under
$\ctor{factor}\;p\;(-)$.
\end{proof}

\begin{lemma}[Evaluation is multiplicative,
  \lean{PrimeTensor.PrimeBag.eval\_mul}]\label{Prime:lem:bag-eval-mul}
$\name{eval}(a \cdot b) = \name{eval}(a) \cdot \name{eval}(b)$ for all
$a, b : \name{PrimeBag}$.
\end{lemma}

\begin{proof}
By structural induction on $a$.

For $a = \ctor{one}$ the left side is $\name{eval}(b)$ by
\eqref{Prime:eq:mul-one-left} and the right side is $1 \cdot \name{eval}(b)$ by
\eqref{Prime:eq:eval-one}; these agree by the left unit law of $\mathbb{N}$.

For $a = \ctor{factor}\;p\;r$ the two sides are, after
\eqref{Prime:eq:mul-factor} and \eqref{Prime:eq:eval-factor},
\[
  p.\name{value} \cdot \name{eval}(r \cdot b)
  \qquad\text{and}\qquad
  \bigl(p.\name{value} \cdot \name{eval}(r)\bigr) \cdot \name{eval}(b).
\]
Rewriting the first by the induction hypothesis makes them
$p.\name{value} \cdot (\name{eval}(r) \cdot \name{eval}(b))$ and
$(p.\name{value} \cdot \name{eval}(r)) \cdot \name{eval}(b)$, equal by
associativity in $\mathbb{N}$.
\end{proof}

\begin{remark}\label{Prime:rem:bag-monoid}
Lemmas~\ref{Prime:lem:bag-one-mul}, \ref{Prime:lem:bag-mul-one} and~\ref{Prime:lem:bag-assoc}
are exactly the monoid laws, and Lemma~\ref{Prime:lem:bag-eval-mul} says
$\name{eval}$ is a monoid homomorphism to $(\mathbb{N}, \cdot, 1)$.  The file
proves the laws but does not register a \lean{Monoid} instance; only
\lean{One} and \lean{Mul} are declared.  What is deliberately absent is
commutativity, which is false for $\name{PrimeBag}$: taking $p \neq q$,
\[
  \ctor{factor}\;p\;(\ctor{factor}\;q\;\ctor{one})
  \;\neq\;
  \ctor{factor}\;q\;(\ctor{factor}\;p\;\ctor{one}),
\]
since $\ctor{factor}$ is injective in its first argument.  (This
non-equality is an observation for the reader; the file states neither it nor
any commutativity law for $\name{PrimeBag}$.)
\end{remark}

\subsection{Identifying factor order}

\begin{definition}[Same, \lean{PrimeTensor.PrimeBag.Same}]\label{Prime:def:same}
For $a, b : \name{PrimeBag}$,
\[
  \name{Same}(a, b) \defeq \bigl(\name{eval}(a) = \name{eval}(b)\bigr).
\]
\end{definition}

$\name{Same}$ is thus, by definition, the kernel of $\name{eval}$: two bags
are identified exactly when they denote the same natural number.

\begin{lemma}[\lean{same\_refl}, \lean{same\_symm}, \lean{same\_trans}]
\label{Prime:lem:same-equiv}
$\name{Same}$ is reflexive, symmetric and transitive.
\end{lemma}

\begin{proof}
Unfolded, each statement is the corresponding property of equality on
$\mathbb{N}$, so the three proofs are $\ctor{rfl}$, $\name{Eq.symm}$ and
$\name{Eq.trans}$ respectively.  They carry \lean{@[refl]}, \lean{@[symm]} and
\lean{@[trans]} so that the relational tactics can use them directly.
\end{proof}

\begin{definition}[The setoid, \lean{PrimeTensor.PrimeBag.setoid}]
\label{Prime:def:setoid}
$\name{setoid} : \name{Setoid}(\name{PrimeBag})$ is $\name{Same}$ together
with the equivalence proof assembled from Lemma~\ref{Prime:lem:same-equiv}.
\end{definition}

\begin{lemma}[Concatenation commutes up to $\name{Same}$,
  \lean{PrimeTensor.PrimeBag.mul\_comm\_same}]\label{Prime:lem:mul-comm-same}
$\name{Same}(a \cdot b,\; b \cdot a)$ for all $a, b : \name{PrimeBag}$.
\end{lemma}

\begin{proof}
Unfolding Definition~\ref{Prime:def:same}, the goal is
$\name{eval}(a \cdot b) = \name{eval}(b \cdot a)$.  Two applications of
Lemma~\ref{Prime:lem:bag-eval-mul} reduce it to
$\name{eval}(a) \cdot \name{eval}(b) = \name{eval}(b) \cdot \name{eval}(a)$,
which is commutativity of multiplication on $\mathbb{N}$.
\end{proof}

\begin{lemma}[$\name{Same}$ is a congruence,
  \lean{PrimeTensor.PrimeBag.mul\_same}]\label{Prime:lem:mul-same}
If $\name{Same}(a_1, a_2)$ and $\name{Same}(b_1, b_2)$ then
$\name{Same}(a_1 \cdot b_1,\; a_2 \cdot b_2)$.
\end{lemma}

\begin{proof}
By Lemma~\ref{Prime:lem:bag-eval-mul} the goal unfolds to
$\name{eval}(a_1) \cdot \name{eval}(b_1)
  = \name{eval}(a_2) \cdot \name{eval}(b_2)$,
and rewriting with the two hypotheses closes it.
\end{proof}

\begin{remark}\label{Prime:rem:cheap-congruence}
Both proofs are short for the same structural reason: $\name{Same}$ was
defined as a property of $\name{eval}$ alone, and $\name{eval}$ was already
known to be multiplicative.  Every fact about the relation is therefore
imported from $\mathbb{N}$ rather than proved about factor chains.  Had
$\name{Same}$ instead been defined as ``is a permutation of'', both lemmas
would have needed genuine combinatorial work.
\end{remark}

\subsection{Prime multisets}

\begin{definition}[\lean{PrimeTensor.PrimeMultiset}]\label{Prime:def:primemultiset}
\[
  \name{PrimeMultiset} \defeq
    \name{Quotient}\bigl(\name{PrimeBag.setoid}\bigr),
\]
declared as an \lean{abbrev}, hence reducible.  Write
$\name{ofBag}(b)$ for the class of $b$, and $\name{one} \defeq
\name{ofBag}(1)$.
\end{definition}

\begin{definition}[Induced operations]\label{Prime:def:quot-ops}
The declarations \lean{PrimeMultiset.eval} and \lean{PrimeMultiset.mul}:
\begin{itemize}[leftmargin=1.5em]
  \item $\name{eval} : \name{PrimeMultiset} \to \mathbb{N}$ is the
    bag-level evaluation of Definition~\ref{Prime:def:bag-eval}, lifted
    through the quotient;
  \item $\name{mul}$ is $(x, y) \mapsto \name{ofBag}(x \cdot y)$ lifted
    through the quotient in both arguments.
\end{itemize}
The instances \lean{One PrimeMultiset} and \lean{Mul PrimeMultiset} give
these the usual notation.
\end{definition}

Each lift carries a well-definedness obligation, and the two are discharged
differently:

\begin{itemize}[leftmargin=1.5em]
  \item For $\name{eval}$ the obligation is that
    $\name{Same}(a, b)$ implies $\name{eval}(a) = \name{eval}(b)$.  By
    Definition~\ref{Prime:def:same} that implication is the identity: the
    hypothesis \emph{is} the conclusion.  This is the payoff of quotienting
    by the kernel of the very map one wants to lift.
  \item For $\name{mul}$ the obligation is that $\name{Same}$ is respected in
    both arguments, which is Lemma~\ref{Prime:lem:mul-same}, followed by
    $\name{Quotient.sound}$ to turn the resulting $\name{Same}$ back into an
    equality of classes.
\end{itemize}

\begin{lemma}[\lean{eval\_ofBag}, \lean{eval\_one}, \lean{eval\_mul}]
\label{Prime:lem:quot-eval}
For $b : \name{PrimeBag}$ and $a, b' : \name{PrimeMultiset}$,
\[
  \name{eval}(\name{ofBag}(b)) = \name{eval}(b),
  \qquad
  \name{eval}(1) = 1,
  \qquad
  \name{eval}(a \cdot b') = \name{eval}(a) \cdot \name{eval}(b').
\]
\end{lemma}

\begin{proof}
The first holds by $\ctor{rfl}$: a lift applied to a representative reduces
definitionally to the underlying function.  The second is the first at
$b = \ctor{one}$ together with \eqref{Prime:eq:eval-one}.  The third is proved by
$\name{Quotient.inductionOn}_2$ --- it suffices to check the identity on
representatives $x, y$, where it becomes
$\name{eval}(x \cdot y) = \name{eval}(x) \cdot \name{eval}(y)$, which is
Lemma~\ref{Prime:lem:bag-eval-mul}.  All three are tagged \lean{@[simp]}.
\end{proof}

\begin{proposition}[Commutativity, \lean{PrimeTensor.PrimeMultiset.mul\_comm}]
\label{Prime:prop:quot-comm}
$a \cdot b = b \cdot a$ for all $a, b : \name{PrimeMultiset}$ --- an equality
of elements, not merely of represented values.
\end{proposition}

\begin{proof}
By $\name{Quotient.inductionOn}_2$ it suffices to prove it for classes of
representatives $x, y$, where the goal is
$\name{ofBag}(x \cdot y) = \name{ofBag}(y \cdot x)$.  By
$\name{Quotient.sound}$ this follows from $\name{Same}(x \cdot y, y \cdot x)$,
which is Lemma~\ref{Prime:lem:mul-comm-same}.
\end{proof}

This is the point of the construction.  Concatenation of chains is not
commutative (Remark~\ref{Prime:rem:bag-monoid}); the quotient makes the induced
operation commutative on the nose, so downstream files may rewrite with
$\name{mul\_comm}$ as an ordinary equation and never carry a
setoid or a ``modulo reordering'' side condition.

\begin{remark}[What the quotient is, and what unique factorization is for]
\label{Prime:rem:ufd}
It is worth being exact about the roles of the two arithmetic facts here,
because the file's header mentions unique factorization while its proofs do
not use it.

Since $\name{Same}$ is the kernel of $\name{eval}$
(Definition~\ref{Prime:def:same}), the induced map
\[
  \name{eval} : \name{PrimeMultiset} \longrightarrow \mathbb{N}
\]
is injective.  By Lemma~\ref{Prime:lem:one-le-eval} its image lies in
$\mathbb{Z}_{\geq 1}$, and since every $n \geq 1$ is a product of primes the
image is all of $\mathbb{Z}_{\geq 1}$.  So $\name{PrimeMultiset}$ is
isomorphic, as a commutative monoid, to $(\mathbb{Z}_{\geq 1}, \cdot, 1)$.

That reorderings are identified by the quotient needs only commutativity of
$\mathbb{N}$, which is what Lemma~\ref{Prime:lem:mul-comm-same} uses.  Unique
factorization is what gives the converse: that \emph{only} reorderings are
identified, so that a class really is the multiset of prime factors and not a
coarser object.  That converse justifies the name $\name{PrimeMultiset}$, but
it is not proved in this file and no result here depends on it.
\end{remark}

\subsection{Concrete atoms and a worked value}

\begin{definition}[\lean{primeTwo}, \lean{primeThree}, \lean{primeFive},
  \lean{twelveBag}, \lean{twelve}]\label{Prime:def:concrete}
$\name{primeTwo}$, $\name{primeThree}$ and $\name{primeFive}$ are the atoms
$2$, $3$ and $5$, each paired with a primality proof discharged by
\lean{norm\_num}.  Further,
\[
  \name{twelveBag} \defeq
    \ctor{factor}\;2\;\bigl(\ctor{factor}\;2\;
      (\ctor{factor}\;3\;\ctor{one})\bigr),
  \qquad
  \name{twelve} \defeq \name{ofBag}(\name{twelveBag}).
\]
\end{definition}

Unfolding Definition~\ref{Prime:def:bag-eval} three times,
\[
  \name{eval}(\name{twelveBag}) = 2 \cdot (2 \cdot (3 \cdot 1)) = 12,
\]
and $\name{primeFive}$ is declared but not used elsewhere in the file.  Note
that $\name{twelveBag}$ fixes an order for the factors of $12$ while
$\name{twelve}$ does not: any of the three arrangements of $2, 2, 3$ gives a
bag with the same image under $\name{ofBag}$, by
Proposition~\ref{Prime:prop:quot-comm}.

\subsection*{Summary of the correspondence}

\begin{center}
\small
\begin{tabular}{@{}lll@{}}
\hline
Lean declaration & Kind & Here \\
\hline
\lean{PrimeTensor.Prime}              & structure  & Def.~\ref{Prime:def:prime} \\
\lean{Prime.one\_lt}                  & simp thm   & Lem.~\ref{Prime:lem:one-lt} \\
\lean{PrimeTensor.PrimeBag}           & inductive  & Def.~\ref{Prime:def:primebag} \\
\lean{PrimeBag.eval}                  & definition & Def.~\ref{Prime:def:bag-eval} \\
\lean{PrimeBag.eval\_one}             & simp thm   & Lem.~\ref{Prime:lem:eval-eqs} \\
\lean{PrimeBag.eval\_factor}          & simp thm   & Lem.~\ref{Prime:lem:eval-eqs} \\
\lean{PrimeBag.one\_le\_eval}         & theorem    & Lem.~\ref{Prime:lem:one-le-eval} \\
\lean{PrimeBag.mul}                   & definition & Def.~\ref{Prime:def:bag-mul} \\
\lean{One PrimeBag}, \lean{Mul PrimeBag} & instances & Def.~\ref{Prime:def:bag-mul} \\
\lean{PrimeBag.one\_mul}              & simp thm   & Lem.~\ref{Prime:lem:bag-one-mul} \\
\lean{PrimeBag.mul\_one}              & simp thm   & Lem.~\ref{Prime:lem:bag-mul-one} \\
\lean{PrimeBag.mul\_assoc}            & theorem    & Lem.~\ref{Prime:lem:bag-assoc} \\
\lean{PrimeBag.eval\_mul}             & simp thm   & Lem.~\ref{Prime:lem:bag-eval-mul} \\
\lean{PrimeBag.Same}                  & definition & Def.~\ref{Prime:def:same} \\
\lean{PrimeBag.same\_refl}            & refl thm   & Lem.~\ref{Prime:lem:same-equiv} \\
\lean{PrimeBag.same\_symm}            & symm thm   & Lem.~\ref{Prime:lem:same-equiv} \\
\lean{PrimeBag.same\_trans}           & trans thm  & Lem.~\ref{Prime:lem:same-equiv} \\
\lean{PrimeBag.setoid}                & definition & Def.~\ref{Prime:def:setoid} \\
\lean{PrimeBag.mul\_comm\_same}       & theorem    & Lem.~\ref{Prime:lem:mul-comm-same} \\
\lean{PrimeBag.mul\_same}             & theorem    & Lem.~\ref{Prime:lem:mul-same} \\
\lean{PrimeTensor.PrimeMultiset}      & abbrev     & Def.~\ref{Prime:def:primemultiset} \\
\lean{PrimeMultiset.ofBag}            & definition & Def.~\ref{Prime:def:primemultiset} \\
\lean{PrimeMultiset.one}              & definition & Def.~\ref{Prime:def:primemultiset} \\
\lean{PrimeMultiset.eval}             & definition & Def.~\ref{Prime:def:quot-ops} \\
\lean{PrimeMultiset.mul}              & definition & Def.~\ref{Prime:def:quot-ops} \\
\lean{One PrimeMultiset}, \lean{Mul PrimeMultiset} & instances & Def.~\ref{Prime:def:quot-ops} \\
\lean{PrimeMultiset.eval\_ofBag}      & simp thm   & Lem.~\ref{Prime:lem:quot-eval} \\
\lean{PrimeMultiset.eval\_one}        & simp thm   & Lem.~\ref{Prime:lem:quot-eval} \\
\lean{PrimeMultiset.eval\_mul}        & simp thm   & Lem.~\ref{Prime:lem:quot-eval} \\
\lean{PrimeMultiset.mul\_comm}        & theorem    & Prop.~\ref{Prime:prop:quot-comm} \\
\lean{primeTwo}, \lean{primeThree}, \lean{primeFive} & definitions & Def.~\ref{Prime:def:concrete} \\
\lean{twelveBag}, \lean{twelve}       & definitions & Def.~\ref{Prime:def:concrete} \\
\hline
\end{tabular}
\end{center}

\noindent
Every declaration of \texttt{PrimeTensor/Prime.lean} appears above.  The file
contains no \lean{sorry}, no axiom, and no unproved named proposition.  The
one statement in this section that is not backed by a declaration in the file
is Remark~\ref{Prime:rem:ufd}'s identification of $\name{PrimeMultiset}$ with
$(\mathbb{Z}_{\geq 1}, \cdot, 1)$, which is stated for the reader.
